\section{Related Work}
\label{sec:relatedWork}
This section provides an overview of previous work on representation learning, retrieval-augmented learning, and generative models.
\label{sec:litreview}
\subsection{Representation Learning}
Deep learning models have achieved superior performance across a wide range of tasks when trained on large collections of well-labeled data. However, the reliance on supervision is often constrained by the cost and difficulty of acquiring large-scale labeled datasets. Representation Learning (RL) provides an effective solution to this challenge. The central idea is to leverage abundant unlabeled or weakly labeled data to learn general-purpose representations that can then be fine-tuned with comparatively few labeled examples.

Contrastive learning approaches, such as SimCLR~\cite{chen2020simple}, demonstrated that powerful visual representations can be obtained by encouraging consistency between augmented views of the same image while distinguishing them from other images. Similarly, Masked Autoencoders (MAE)~\cite{he2022masked} revealed that randomly masking a large portion of an image and training a model to reconstruct the missing content from the visible context can yield highly effective representations. More recently, pre-training using multimodal data, eg: image/text, image/geolocation, audio/satellite-image) has emerged as a powerful RL strategy. 

\subsubsection{Vision-Language Pre-training}
Vision-Language (VL) models have shown great promise in their ability to model complex relationships between the vision and text space. ConVIRT~\cite{zhang2022contrastive} and VirTex~\cite{desai2021virtex} both introduced methods that used image-text pairs to learn rich visual representations. CLIP~\cite{radford2021learning} demonstrated the results of VL pretraining on a large-scale dataset (400M pairs) and validated the efficacy of large-scale VL pretraining for several downstream tasks. Florence~\cite{yuan2021florence} and ALIGN~\cite{jia2021scaling} further increased the scale of data by training on 900M and 1.8B pairs, respectively. Other works~\cite{li2021align,yao2021filip,yang2022unified,yang2022vision} have since focused on learning a better VL embedding space. With the existence of these powerful pretrained VL models, many researchers have utilized their embedding spaces to solve specific downstream tasks. CLIPCap~\cite{mokady2021clipcap} and~\cite{cho2022fine} used CLIP space to generate image captions. Models like~\cite{ramesh2022hierarchical,nichol2021glide,wang2022clip} utilized the CLIP space for text-to-image generation. 

Recently, there has been work in creating image-text datasets of overhead images and captions for VL pertaining of geospatial models. ChatEarthNet~\cite{yuan2024chatearthnet} created a dataset with paired image and text captions using ChatGPT. Similarly, SkyScript~\cite{wang2023skyscript} used OpenStreetMap (OSM) data to create overhead image and text-paired datasets for VL pretraining. However, both approaches only gather coarse textual information for a given location using either the low-resolution visual cues from an overhead image or fixed preset tags from OSM data.

Our approach fundamentally differs from prior work as we try to capture the more intricate concepts occurring at the ground-level of a given location by leveraging the corresponding crowdsourced images. Utilizing the ground-level images uploaded from a location allows us to access extremely fine-grained information of that location that is not distillable solely from satellite images or OSM data.  

\subsubsection{Learning Representations for Geolocation}
Learning effective representations for geolocation is valuable for diverse downstream geospatial tasks, such as geolocalization~\cite{vivanco2024geoclip}, canopy height prediction~\cite{lang2023high}, fine-grained species classification~\cite{Aodha_2019_ICCV,sastry2024birdsat}, species distribution modeling~\cite{cole2023spatial, sastry2024birdsat}, geographical question answering \cite{mai2020se}, soundscape mapping~\cite{khanal2024psm}, etc. Geolocation in these tasks can be represented using either non-parametric~\cite{gaolearning, russwurm2023geographic} or parametric~\cite{mai2023csp,vivanco2024geoclip,klemmer2023satclip} methods.
Ru{\ss}wurm \textit{et al.}~\cite{russwurm2023geographic} represents geolocation using spherical harmonic basis functions, and Aodha \textit{et al.}~\cite{Aodha_2019_ICCV} uses a sine and cosine encoding of geolocation that removes discontinuity at the dateline.

In recent years, learned representations for geolocation, typically parameterized by neural networks, have been popularized. Some of these methods learn geolocation representations by using contrastive pre-training strategies that align location information with visual information from either co-located overhead~\cite{klemmer2023satclip,mai2023csp} or ground-level images~\cite{mai2023csp,vivanco2024geoclip}. The location encoders trained by these methods have demonstrated utility across various downstream geospatial tasks. We use the representations from these location encoders as baselines for our model.

\subsubsection{Multimodal Contrastive Learning}
Contrastive learning is an efficient and scalable pre-training strategy for learning a shared embedding space across multiple modalities~\cite{radford2021learning,yang2022vision,yaofilip,li2021align,girdhar2023imagebind,elizalde2023clap,Dhakal_2024_CVPR,khanal2024psm,huynh2024contrastive,zheng2023exif}. As demonstrated by CLIP~\cite{radford2021learning} and more recent works~\cite{li2021align,yaofilip,yang2022vision} in vision-language learning, the shared embedding space can be utilized for zero-shot classification and leveraged across a variety of downstream tasks. Contrastively trained multimodal embedding spaces between geolocation and satellite imagery~\cite{klemmer2023satclip, mai2023csp} or ground-level imagery~\cite{vivanco2024geoclip}, have also demonstrated impressive transferability of both location and visual representations to various geospatial tasks.

However, despite the success of contrastive pre-training, recent work highlights the limitations of multimodal contrastive learning. One such study~\cite{liang2023factorized} shows that contrastively trained representations capture the shared information between the two modalities and thus only work best under certain assumptions. Similarly, another work~\cite{gupta2022understanding} demonstrates that the final layers of encoders in the contrastive learning framework preserve only the information necessary for solving the alignment objective, discarding other modality-specific details. We argue that some of the underlying assumptions of multimodal contrastive training do not hold true in the image-location setting. Therefore, the resulting location representations from these methods are sub-optimal for downstream tasks. To address this issue, we propose a method to approximate the image-specific information for a given location.


\subsection{Retrieval Augmented Methods}
There has been a recent surge in retrieval-augmented generation (RAG) methods, primarily focused on generative tasks~\cite{chenre, blattmann2022retrieval, kirstain2023x, lewis2020retrieval, borgeaud2022improving, siriwardhana2023improving, seoretrieval, koizumi2020audio, wangretrieval}. These methods generally consist of three main components: a retriever, an external database, and a generator. For a given query, the retriever selects a set of documents from the database and provides them, along with the original query, to the generator to produce the desired output. RAG has been successfully employed to improve text generation for large language models (LLMs)~\cite{lewis2020retrieval, borgeaud2022improving, siriwardhana2023improving}, text-to-image generation~\cite{chenre, blattmann2022retrieval, kirstain2023x}, text-to-3D generation~\cite{seoretrieval}, audio captioning~\cite{koizumi2020audio}, and more. Within the RAG framework, one could also replace the generator with another task-specific component, such as a classifier~\cite{long2022retrieval}, to enhance performance by leveraging the rich information from the retrieved documents. Inspired by these retrieval-augmented methods, we design our framework, RANGE, which retrieves rich visual information from our database for a given geolocation query and utilizes it for various downstream geospatial tasks.

\subsection{Deepfake Generation and Detection}
There has been significant progress in the field of image generative models over the last five years. Earlier approaches were dominated by Generative Adversarial Networks~\cite{goodfellow2020generative}. Recently, however, Diffusion Models~\cite{ho2020denoising} have emerged as the state-of-the-art framework for generative modeling. Several works have demonstrated the ability of diffusion models to generate highly photorealistic images using various input conditions, including text, bounding boxes, and edge maps~\cite{rombach2022high, saharia2022photorealistic, peebles2023scalable,zhang2023adding}. While there has been a tremendous amount of work on general-purpose diffusion models, several researchers have also focused on training specialized diffusion models to generate images in specific data domains such as satellite images~\cite{khanna2023diffusionsat, sastry2024geosynth, tang2024crs}, medical images~\cite{khader2023denoising, yoon2023sadm}, spectrograms~\cite{majumder2024tango, ghosal2023text}, etc. In our work, we aim to understand the efficacy and limitations of specialized generative models for satellite images using powerful pre-trained models.    

As generative models become increasingly powerful, there is a growing need for robust deepfake detection methods. However, several challenges need to be addressed to get deep-fake detectors with practical applications. Many methods tend to overfit to models present in their training set, learning trivial solutions. As deepfake models typically require generated images for training, they also lag behind the introduction of new generative models, which can be fatal in mission-critical scenarios such as detecting deep-fake satellite images. We aim to propose deep-fake detectors that generalize well outside the training set models, under limited training data settings. 